\chapter{Introduction}
\label{sec1:introduction}

The rapid global emergence of antimicrobial resistance (AMR) represents one of the most critical threats that modern 
medicine is facing. As conventional antibiotics lose efficiency against multidrug resistant bacterial pathogens, alternative 
therapeutic strategies have gained widespread attention. Among these is phage therapy, it is the targeted clinical application of 
bacteriophages (viruses that specifically infect and destroy bacteria) and has emerged as a highly promising solution 
\cite{dejonge2019molecular}.

Bacteriophages are the most abundant biological entities in the global biosphere, with an estimated population exceeding 
$10^{31}$ particles \cite{dejonge2019molecular}. One of the main determinants of a phage's therapeutic viability is its host range 
(the spectrum of bacterial strains it can successfully infect), so matching a bacterial infection with an effective phage 
requires knowing precisely which bacteria each phage can infect and which it cannot. Section~\ref{sec2.1: host prediction problem} reviews the biology of phages and of the host-prediction problem in detail.

% Molecular recognition is governed by highly specialized 
% interactions between viral structural proteins (such as tail fibers, lysins etc.) and bacterial cell-wall receptors, which 
% suggests that infection should be highly specific. However, in practice, phage host specificity is not absolute, the range spans 
% from single bacteria type to multiple species or genera simultaneously. Yet, having a narrow host range is a major clinical asset, as 
% it allows therapeutic phages to kill target pathogens without harming the patient, but it also introduces a diagnostic 
% challenge, matching a bacterial infection with an effective phage requires knowing precisely which bacteria each phage can infect.

Traditional laboratory screening methods, such as plaque assays, are labor intensive, costly, and scale poorly against 
the millions of uncharacterized viral genomes generated by modern high throughput metagenomic sequencing. This necessitates 
the development of fast and scalable computational tools that take genomic sequences as input, translate them into protein 
level representations and then predict the host range from these representations directly without requiring any 
laboratory validation \cite{coclet2021global}.

However, state of the art computational prediction tools suffer from a fundamental limitation. They produce point 
predictions or uncalibrated scalar scores in $[0,1]$ without any finite sample uncertainty quantification. In high stakes 
medical applications, point predictions offer no guarantee of reliability; a confidence score reflects only a relative ranking of 
candidates, not a calibrated probability that the predicted host is correct. The clinical risk runs in both directions, failing to 
recognize that a phage kills the infecting bacterium can mean missing a life saving treatment, while overlooking that the same 
phage may also attack beneficial bacteria can cause further damage. Therefore, sound therapeutic decisions require calibrated 
knowledge of which phages infect which bacteria, and equally of which they do not; point predictions alone are clinically 
insufficient.


The main objective of this work is to resolve this limitation by constructing an uncertainty aware statistical framework 
for phage-host prediction using conformal prediction \cite{vovk2005algorithmic, angelopoulos2023conformal}, which is a 
distribution free method that transforms point scoring models into set valued predictions $C(X) \subseteq \mathcal{H}$. 
More importantly, these prediction sets are mathematically guaranteed to contain the true host label with a user specified 
confidence level $1-\alpha$ (e.g., $90\%$), which holds strictly in finite samples under just the assumption of data exchangeability.

% This guarantee, however, is stated for the classical setting of a single predictor that assigns each example one real 
% valued nonconformity score. Phage-host prediction does not fit this template: each candidate host is assessed by 
% several specialized scoring functions (one per protein family, e.g., capsid, lysin, tail fiber), so a phage is 
% described by an entire vector of scores rather than by a single number. Adapting conformal prediction to this multiple 
% predictor regime therefore requires deciding how the score vector enters the validity guarantee. One option is to c
% ollapse the vector into a scalar, for instance by averaging, and apply vanilla split conformal prediction, but this 
% discards the complementary information carried by the individual functions. The alternative is to keep the full score 
% vector and define a joint acceptance region in score space, calibrated so that the true label's score vector falls inside 
% it with probability at least $1-\alpha$. This second route preserves far more information, but it raises an explicitly 
% geometric question: what shape should this acceptance region take? Answering this question is precisely the purpose of the 
% Conformal Set Aggregation (CSA) framework, to which we now turn.

Moreover, standard conformal prediction relies on a single score to evaluate each candidate label, whereas our pipeline evaluates 
phage host compatibility across $K$ distinct functional protein categories simultaneously. Simply averaging these scores into one 
value loses useful information, while calibrating each score separately breaks the coverage guarantee; the multi dimensional 
conformal methods that operate directly on score vectors are introduced in Section~\ref{sec2.4: CSA Method}.

In this work we build on the Conformal Set Aggregation (CSA) framework \cite{ochoarivera2024conformal}. We address 
three major challenges that arise when applying CSA to classification datasets:

\begin{enumerate}[label=(\roman*)]
    \item \textbf{Non-convex geometry in classification space:} We demonstrate why traditional convex envelopes, developed 
           primarily for continuous regression spaces, are inefficient in classification settings, since the class conditional score 
           vectors occupy disjoint, bounded clusters in $[0,1]^K$. So, we propose non-convex (Radial and Strip) and Collapsed 1D 
           envelope geometries that adapt to cluster topologies, and eliminate unpopulated empty space, which helps us reduce 
           prediction set sizes.
    
    \item \textbf{Closed form analytical calibration ($\tau$):} We derive exact, closed form analytical expressions for the 
           per point inclusion scaling function $\tau(\mathbf{s})$ across all envelope geometries. This eliminates the 
           computationally expensive numerical binary search required by standard CSA, which reduces the calibration time 
           complexity from $O(|S_2| \cdot N_{\text{iter}})$ to a single forward pass and $O(|S_2| \log |S_2|)$ sort.
    
%     \item \textbf{Class Conditional Validity:} Global calibration can yield under coverage for complex host classes and over 
%            coverage for simpler ones. To prevent this we incorporate per class quantile calibration ($\hat{t}_c$). This guarantees 
%            exact $1-\alpha$ coverage conditionally for every host label individually rather than just on average.
    
    \item  \textbf{Structural missingness handling:} Phage score vectors frequently contain entries that are missing 
            by design, since a phage may lack any protein annotated for a given functional category. We develop the 
            envelope constructions and nonconformity transformations to operate natively on these partially observed 
            vectors via an explicit missingness mask, without imputation. 
       %      This treatment is not specific to phage 
       %      biology and applies directly to other multi dimensional conformal prediction problems with structural 
       %      (rather than random) missingnes
\end{enumerate}

We evaluated our methodology systematically on two tasks of increasing complexity. First, 
a controlled binary Gram-type classification task ($N=2$ classes) to validate the geometric envelope constructions, and then on a 
multi class host-level prediction task ($N=54$ bacterial genera) using real-world protein language model embeddings.


The remainder of this report is organized as follows: Chapter~\ref{sec2: background} reviews the biological background 
of bacteriophages, computational state-of-the-art methods, and the mathematical foundations of split conformal prediction and CSA. 
Chapter~\ref{sec3: methodology} formalizes the prediction tasks, develops the non-convex envelopes, the nonconformity 
transformation, and the analytical derivation of $\tau(\mathbf{s})$. Chapter~\ref{sec4: data} details the dataset architecture, 
protein to phage NaN aware mean pooling operators, and cross split protein cluster contamination filtering. Chapter~
\ref{sec5: experiments} presents synthetic experiments and empirical evaluation results for both gram-type and host-level 
classification tasks. Finally, chapter~\ref{sec6: conclusion} summarises our findings, discusses the biological limitations, and 
outlines directions for future research.
%%%%%%%%%%%%%%%%%%%%%%%%%%%%%%%%%%%%%%%%%%%%%%%%%%%%%%%%%%%%%%%%%%%%%%%%%%%%%%%%%%%%%


